%% Elevark fra mattelab.no – laget automatisk av bygg.py. Rediger gjerne til eget bruk.
%% Kompiler med lualatex.
\documentclass[11pt,a4paper]{article}
\usepackage[norsk]{babel}
\usepackage[a4paper,margin=2.2cm,top=2cm,bottom=2.4cm]{geometry}
\usepackage{fontspec}
\usepackage{unicode-math}
\setmainfont{TeX Gyre Pagella}
\setsansfont{TeX Gyre Heros}[Scale=MatchLowercase]
\setmathfont{TeX Gyre Pagella Math}
\usepackage{xcolor}
\definecolor{blekk}{HTML}{16213A}
\definecolor{dempet}{HTML}{566077}
\definecolor{rute}{HTML}{CFC9B8}
\definecolor{m1}{HTML}{2F4AB0}
\definecolor{m2}{HTML}{99266D}
\usepackage[most]{tcolorbox}
\usepackage{qrcode}
\usepackage{needspace}
\usepackage{fancyhdr}
\usepackage[hidelinks]{hyperref}
\pdfvariable suppressoptionalinfo \numexpr 1+2+4+8+16+32+64+128+256+512\relax
\setlength{\parindent}{0pt}
\setlength{\parskip}{0.4em}
\pagestyle{fancy}
\fancyhf{}
\renewcommand{\headrulewidth}{0pt}
\fancyfoot[L]{\sffamily\footnotesize\color{dempet}mattelab.no · Derivasjon av eksponential- og logaritmefunksjoner · elevark uten fasit}
\fancyfoot[R]{\sffamily\footnotesize\color{dempet}Side \thepage}

\newcommand{\fase}[2]{\tcbox[on line, colframe=#2, colback=white, coltext=#2, boxrule=0.6pt, arc=2pt,
  boxsep=0pt, left=3pt, right=3pt, top=1.5pt, bottom=1.5pt, fontupper=\sffamily\bfseries\footnotesize]{#1}}
\newcommand{\steg}[4]{%
  \par\noindent{\color{rute}\rule{\linewidth}{0.4pt}}\par\nopagebreak
  \noindent\begin{minipage}[t]{2.1cm}\vspace{0pt}\fase{#1}{#2}\end{minipage}%
  \begin{minipage}[t]{\dimexpr\linewidth-2.1cm\relax}\vspace{0pt}#3\end{minipage}\par
  \vspace{#4}}
\newcommand{\del}[4]{%
  \par\noindent{\color{rute}\rule{\linewidth}{0.4pt}}\par\nopagebreak
  \noindent\begin{minipage}[t]{1.1cm}\vspace{0pt}\fase{#1}{#2}\end{minipage}%
  \begin{minipage}[t]{\dimexpr\linewidth-1.1cm\relax}\vspace{0pt}#3\end{minipage}\par
  \vspace{#4}}
\newcommand{\oppgave}[2]{\par\Needspace{12\baselineskip}\vspace{0.6em}%
  {\fontsize{14}{17}\selectfont\bfseries{\color{m1}#1}\enspace #2}\par\vspace{0.2em}}

\begin{document}
\begin{minipage}[t]{0.72\linewidth}\vspace{0pt}
{\sffamily\small\color{dempet}Elevark · R1 · mattelab.no}\par\vspace{0.3em}
{\fontsize{24}{28}\selectfont\bfseries Derivasjon av eksponential- og logaritmefunksjoner}\par\vspace{1em}
{\sffamily\small Navn: \rule{4.6cm}{0.4pt}\quad Klasse: \rule{1.5cm}{0.4pt}\quad Dato: \rule{1.9cm}{0.4pt}}
\end{minipage}\hfill
\begin{minipage}[t]{0.22\linewidth}\vspace{0pt}\raggedleft
\qrcode[height=2.3cm]{https://mattelab.no/opplegg/derivasjon-av-eksponential-og-logaritmefunksjoner.html}\par\vspace{0.2em}
{\sffamily\scriptsize\color{dempet}Opplegget}
\end{minipage}

\vspace{1em}
{\sffamily\small\color{dempet}Opplegget ligger på \url{https://mattelab.no/opplegg/derivasjon-av-eksponential-og-logaritmefunksjoner.html}}\par
Oppgave 1 og 3 skal løses uten hjelpemidler. I oppgave 2 kan du bruke lommeregner.

\oppgave{1}{Deriver}
\del{a)}{blekk}{$f(x) = 5e^{2x}$}{2.5cm}
\del{b)}{blekk}{$g(x) = \ln(4x)$}{2.5cm}
\del{c)}{blekk}{$h(x) = 3 \cdot 2^x$. Regn også ut $h'(0)$ med to desimaler.}{2.5cm}

\oppgave{2}{Teen blir kald igjen}
\del{a)}{blekk}{Temperaturen i en kopp te er $T(t) = 20 + 60e^{-0{,}1t}$ grader etter $t$ minutter. Finn $T'(t)$.}{2.5cm}
\del{b)}{blekk}{Regn ut $T'(5)$, og forklar hva tallet betyr.}{2.5cm}
\del{c)}{blekk}{Når synker temperaturen med 1 grad per minutt?}{2.5cm}

\oppgave{3}{Et bunnpunkt}
\del{a)}{blekk}{La $f(x) = x - \ln x$ for $x > 0$. Finn $f'(x)$.}{2.5cm}
\del{b)}{blekk}{Vis at grafen har et bunnpunkt, og finn det.}{2.5cm}

\end{document}
